\section{Conclusion}
\label{sec:conclusion}
This work applied Persistence Mechanics to the most fundamental physical baseline: the quantum vacuum. We tested whether the requirement of resisting entropic decay is sufficient to define the substrate of reality. Implementing the three existential requirements (Finiteness, Conservation, and Causality) through its six-pillar architecture strictly constrained the viable vacuum substrate to the $E_8$ lattice projected onto a causal 4D manifold. 

This purely architectural derivation intrinsically yields the spacetime geometry including the Lorentzian metric signature, a fundamental temporal cycle ($\Delta = 43$), discrete channel depth ($\nu = 16$), embedding dimension ($\sigma = 5$), and boundary topology ($\chi = 2$) as exact mathematical consequences, with no free parameters. Beyond the structural details it is worth highlighting that applying persistence mechanics resulted in the isolation of a single substrate.

The second, decisive test evaluates whether this substrate's geometric impedance $Z_{\text{geo}}$ (a structural index derived from the six-pillar entropic balance) correctly corresponds to the static, fine-structure constant ($\alpha^{-1}(0) = \num{\AlphaInvVal}\ldots$). This value structurally aligns with the kinematic measurement of Morel (2020) at $\AlphaInvMorelSigma\sigma$ and CODATA (2022) at $\AlphaInvCodataSigma\sigma$, while predicting a strict \textit{Quantization Gap} where continuous, high-loop QED extractions must systematically diverge due to the substrate's finite channel capacity. The same six-pillar impedance is mapped consistently to the von Klitzing constant $R_K$ and the elementary charge $e = q_P/\sqrt{Z_{geo}}$, which emerges as a flow constraint enforcing the Schwinger vacuum breakdown limit.

Having established the static baseline of the $E_8$ substrate, it serves as a derivational platform: these extensions provide an explicit roadmap to verify the emergence of General Relativity, the Standard Model Lagrangian, and its remaining dimensionless constants, which can then be compared against experimental values. Ultimately, the definitive proof that $\alpha^{-1}$ is an architectural necessity lies in whether this exact, unadjusted substrate generates the full Standard Model action and remaining constants.

These initial results have striking implications for fundamental physics if they continue to hold. Physical constants are traditionally treated as arbitrary empirical parameters requiring external measurement. The $E_8$-Persistence Theory suggests instead that they are rigid geometric boundaries. The vacuum is not a passive empty space; it is a structural architecture actively solving the problem of its own existence.

If the fundamental vacuum is uniquely determined by the strict constraints of persistence, then by the Selection Principle, any persistent macroscopic system is solving the exact same architectural problem, differing only in substrate material and environmental openness.

\begin{acknowledgments}
The author is an independent researcher and received no external funding for this work. I would like to thank Brian Sheppard for rigorous and constructive feedback.
\end{acknowledgments}